\section{Resultados}
    En la presente sección se exponen los resultados más relevantes de cada 
    imagen del conjunto, destacando los filtros que obtuvieron mejores 
    resultados, así como explicando sus ventajas y desventajas específicas 
    en cada imagen o escenario. 

    \subsection{Imagen 1}
        Esta imagen contiene elementos tanto geométricos como no geométricas 
        con texturas muy variables. Esto significa que los filtros que sean 
        sensibles al ruido, van a detectar muchos no bordes, como le pasa 
        principalmente al filtro de Laplaciano de Gaussiana; esto también 
        ocurre en los filtros de Prewitt y Sobel. Pero este fenómeno se 
        puede mitigar con un adecuado suavizado gaussiano.

        \begin{figure}[H]
            \centering
            \includegraphics[width=0.7\textwidth]{../results/describe_1.jpg}
            \caption{Imagen 1 del conjunto.}
        \end{figure}

        \begin{figure}[H]
            \centering
            \includegraphics[width=0.65\textwidth]{../results/filter_laplacian_gaussian_1.jpg}
            \caption{Filtro Laplaciano de Gaussiana sobre la Imagen 1 sin suavizado.}
        \end{figure}

        Al aplicar un suavizado gaussiano con un kernel de $7\times7$, se pueden 
        mejorar los bordes detectados, haciendo que se vean más definidos o sin 
        tanto ruido debido a las texturas. Donde más se muestra una mejora es en 
        el filtro de Canny, debido a que detecta adecuadamente los bordes sin 
        añadir ruido por la textura de los elementos, dicho así, el combinar la 
        información de diferentes filtros le favorece para incrementar la sensibilidad 
        al momento de detectar un verdadero borde. 

        \begin{figure}[H]
            \centering
            \includegraphics[width=0.65\textwidth]{../results/filter_canny_1_suavizado_7x7.jpg}
            \caption{Filtro de Canny sobre la Imagen 1 suavizada con un kernel $7\times7$.}
        \end{figure}

    \subsection{Imagen 2}
        Esta imagen cuenta con dos peculiaridades: tiene una sobreexposición 
        y vegetación que añade mucha texturas y contraste de colores (áreas 
        oscuras y claras). Por lo tanto, para reducir el efecto del brillo, se 
        le aplicó una función de ganancia-sesgo, esto va a favorece a encontrar 
        algunos de los bordes empleando los filtros adecuados. Sin aplicar un 
        suavizado gaussiano, cualquier filtro detecta parte del ruido presente, 
        donde el filtro Laplaciano de Gaussiana es el más sensible a este efecto, 
        y los filtros de Prewitt, Sobel y Canny manejan de mejor manera el ruido.

        \begin{figure}[H]
            \centering
            \includegraphics[width=0.7\textwidth]{../results/describe_2_enhancement.jpg}
            \caption{Imagen 2 del conjunto.}
        \end{figure}

        \begin{figure}[H]
            \centering
            \includegraphics[width=0.65\textwidth]{../results/filter_prewitt_2.jpg}
            \caption{Filtro Prewitt sobre la Imagen 2 sin suavizado.}
        \end{figure}

        Después de aplicar un suavizado gaussiano con un kernel de $7\times7$, es 
        posible reducir el efecto de la textura de la vegetación y se aprovecha de 
        mejor manera los filtros de Sobel y Prewitt para delimitar los elementos. 
        Con este suavizado, los filtros de Laplace y Laplaciano de Gaussiana ignoran 
        los cambios más sensibles de los colores (o locales).

        \begin{figure}[H]
            \centering
            \includegraphics[width=0.65\textwidth]{../results/filter_sobel_2_suavizado_7x7.jpg}
            \caption{Filtro de Sobel sobre la Imagen 2 suavizada con un kernel $7\times7$.}
        \end{figure}

    \subsection{Imagen 3}
        Esta imagen dominan los elementos geométricos, esto permite que los filtros 
        de Prewitt, Sobel y Canny detecten los bordes adecuadamente, mientras que 
        los otros filtros puede que detecten algo de ruido por las texturas de las 
        maderas. Con el fin de facilitar la detección de bordes, se mejora el brillo 
        de la imagen con la función de ganancia-sesgo. Sin aplicar un suavizado, 
        los filtros mencionados detectan los bordes adecuadamente, pero capturan algo 
        de ruido por la textura de los elementos, siendo esta su mayor desventaja al 
        ser sensibles a los cambios entre píxeles. Mientras que, en este caso, el filtro 
        de Canny se vuelve el más robusto contra el ruido y funciona adecuadamente 
        para ignorar falsos bordes.

        \begin{figure}[H]
            \centering
            \includegraphics[width=0.7\textwidth]{../results/describe_3_enhancement.jpg}
            \caption{Imagen 3 del conjunto.}
        \end{figure}

        \begin{figure}[H]
            \centering
            \includegraphics[width=0.65\textwidth]{../results/filter_canny_3.jpg}
            \caption{Filtro Canny sobre la Imagen 3 sin suavizado.}
        \end{figure}

        Con un suavizado gaussiano con kernel de $3\times3$, se logra que el filtro 
        Laplaciano de Gaussiana ignore ligeramente la textura de los elementos y detecte 
        adecuadamente los bordes, notando así que este filtro es sensible al ruido o 
        cambios locales de colores.

        \begin{figure}[H]
            \centering
            \includegraphics[width=0.65\textwidth]{../results/filter_laplacian_gaussian_3_suavizado_3x3.jpg}
            \caption{Filtro de Laplaciano de Gaussiana sobre la Imagen 3 suavizada con un kernel $3\times3$.}
        \end{figure}
        
    \subsection{Imagen 4}
        Esta imagen está dominada por elementos geométricos (rectas), esto 
        permite que todos los filtros detecten los bordes más relevantes que 
        constituyen a la imagen. Aunque cuenta con zonas oscuras y claras, éstas 
        no afectan a la detección de bordes, debido a que no obstruyen a las 
        diferencias entre tonalidades (entre un ventana y  una pared). En el 
        caso de los filtros basados en la segunda derivada, tienen la desventaja 
        de captura algo del ruido de las texturas, mientras que Prewitt y Sobel 
        delimitan adecuadamente los bordes de los elementos.

        \begin{figure}[H]
            \centering
            \includegraphics[width=0.7\textwidth]{../results/describe_4.jpg}
            \caption{Imagen 4 del conjunto.}
        \end{figure}

        \begin{figure}[H]
            \centering
            \includegraphics[width=0.65\textwidth]{../results/filter_laplacian_gaussian_4.jpg}
            \caption{Filtro Laplaciano de Gaussiana sobre la Imagen 4 sin suavizado.}
        \end{figure}

        Reduciendo el ruido con un suavizado gaussiano con kernel de $7\times7$, 
        se logra que el filtro de Canny detecte los bordes de los elementos principales 
        que se aprecian con un buen brillo, es decir, aquellos que están en la 
        zonas claras sin que las texturas creen falsos bordes. 

        \begin{figure}[H]
            \centering
            \includegraphics[width=0.65\textwidth]{../results/filter_canny_4_suavizado_7x7.jpg}
            \caption{Filtro de Canny sobre la Imagen 4 suavizada con un kernel $7\times7$.}
        \end{figure}

    \subsection{Imagen 5}
        Aunque en esta imagen los elementos principales tiene una buena 
        diferenciación de tonalidades, podemos notar que también cuentan 
        con muchas texturas que harán que cualquier filtro las vea como 
        bordes. Por ello, sin aplicar el suavizado cualquier filtro detecta 
        esos cambios menores entre colores, lo que añade ruido a los bordes 
        de forma innecesaria. 

        \begin{figure}[H]
            \centering
            \includegraphics[width=0.7\textwidth]{../results/describe_7.jpg}
            \caption{Imagen 5 del conjunto.}
        \end{figure}

        \begin{figure}[H]
            \centering
            \includegraphics[width=0.65\textwidth]{../results/filter_prewitt_7.jpg}
            \caption{Filtro Prewitt sobre la Imagen 5 sin suavizado.}
        \end{figure}

        Aplicando un suavizado gaussiano con kernel de $7\times7$, permite 
        detectar bordes correctamente sin el ruido añadido por las texturas 
        de los diferentes elementos presentes en la imagen. En este caso, 
        los filtros de Prewitt, Sobel y Canny detectan los bordes relevantes 
        de los objetos sin dificultades, mientras que los filtros de Laplace y 
        Laplaciano de Gaussiana no encuentran todos los bordes, tal vez debido 
        a que el suavizado difuminó demasiado las tonalidades.

        \begin{figure}[H]
            \centering
            \includegraphics[width=0.65\textwidth]{../results/filter_canny_7_suavizado_7x7.jpg}
            \caption{Filtro de Canny sobre la Imagen 5 suavizada con un kernel $7\times7$.}
        \end{figure}